% ============================================================
% IAET-RR — Apresentação Beamer
% Indicador de Atividade Econômica Trimestral de Roraima
% SEPLAN/RR — DIEAS/CGEES
% Compile no Overleaf (pdflatex)
% ============================================================

\documentclass[aspectratio=169, 11pt]{beamer}

% --- Tema e cores -------------------------------------------
\usetheme{Madrid}
\usecolortheme{default}

\definecolor{seplanblue}{RGB}{20, 52, 100}
\definecolor{seplanblue2}{RGB}{32, 78, 145}
\definecolor{seplangold}{RGB}{200, 145, 20}
\definecolor{seplanlightblue}{RGB}{220, 232, 248}
\definecolor{seplangreen}{RGB}{34, 130, 64}
\definecolor{seplanred}{RGB}{180, 40, 40}
\definecolor{seplantext}{RGB}{30, 30, 50}

\setbeamercolor{palette primary}{bg=seplanblue, fg=white}
\setbeamercolor{palette secondary}{bg=seplanblue2, fg=white}
\setbeamercolor{palette tertiary}{bg=seplangold, fg=white}
\setbeamercolor{palette quaternary}{bg=seplanblue, fg=white}
\setbeamercolor{structure}{fg=seplanblue}
\setbeamercolor{section in toc}{fg=seplanblue}
\setbeamercolor{frametitle}{bg=seplanblue, fg=white}
\setbeamercolor{title}{bg=seplanblue, fg=white}
\setbeamercolor{block title}{bg=seplanblue2, fg=white}
\setbeamercolor{block body}{bg=seplanlightblue, fg=seplantext}
\setbeamercolor{block title alerted}{bg=seplanred, fg=white}
\setbeamercolor{block body alerted}{bg=red!5, fg=seplantext}
\setbeamercolor{block title example}{bg=seplangreen, fg=white}
\setbeamercolor{block body example}{bg=green!5, fg=seplantext}
\setbeamercolor{itemize item}{fg=seplanblue}
\setbeamercolor{itemize subitem}{fg=seplangold}

\setbeamerfont{frametitle}{size=\large, series=\bfseries}
\setbeamerfont{title}{size=\LARGE, series=\bfseries}

% --- Pacotes ------------------------------------------------
\usepackage[utf8]{inputenc}
\usepackage[T1]{fontenc}
\usepackage[brazil]{babel}
\usepackage{pgfplots}
\pgfplotsset{compat=1.18}
\usepackage{pgf-pie}
\usepackage{booktabs}
\usepackage{multirow}
\usepackage{colortbl}
\usepackage{array}
\usepackage{tikz}
\usetikzlibrary{arrows.meta, positioning, shapes.geometric, fit, backgrounds}
\usepackage{fontawesome5}
\usepackage{xcolor}
\usepackage{soul}
\usepackage{amsmath}
\usepackage{lmodern}

% --- Remover navegação --------------------------------------
\setbeamertemplate{navigation symbols}{}
\setbeamertemplate{footline}{%
  \leavevmode%
  \hbox{%
    \begin{beamercolorbox}[wd=.35\paperwidth,ht=2.5ex,dp=1.5ex,left,leftskip=1em]{palette primary}%
      \usebeamerfont{author in head/foot}\insertshortauthor
    \end{beamercolorbox}%
    \begin{beamercolorbox}[wd=.45\paperwidth,ht=2.5ex,dp=1.5ex,center]{palette secondary}%
      \usebeamerfont{title in head/foot}\insertshorttitle
    \end{beamercolorbox}%
    \begin{beamercolorbox}[wd=.20\paperwidth,ht=2.5ex,dp=1.5ex,right,rightskip=1em]{palette primary}%
      \usebeamerfont{date in head/foot}\insertframenumber{} / \inserttotalframenumber
    \end{beamercolorbox}}%
}

% --- Metadados ----------------------------------------------
\title[IAET-RR]{Indicador de Atividade Econômica\\Trimestral de Roraima}
\subtitle{IAET-RR — Inteligência Econômica em Tempo Real para o Estado}
\author[Yuri Cesar de Lima e Silva]{Yuri Cesar de Lima e Silva\\
  {\small Chefe da DIEAS | Coordenador da Equipe PIB-RR}}
\institute[SEPLAN-RR]{%
  Secretaria de Estado do Planejamento e Desenvolvimento de Roraima\\
  Coordenação-Geral de Estudos Econômicos e Sociais — CGEES}
\date{Abril de 2026}

% ============================================================
\begin{document}
% ============================================================

% --- Slide de título ----------------------------------------
{
\setbeamertemplate{footline}{}
\begin{frame}[plain]
  \titlepage
  \vspace{-0.5cm}
  \begin{center}
    \small\color{seplanblue}
    \textit{Apresentação para a Coord. da CGEES e o Secretário da SEPLAN/RR}
  \end{center}
\end{frame}
}

% --- Sumário ------------------------------------------------
\begin{frame}{Roteiro da Apresentação}
  \tableofcontents[hideallsubsections]
\end{frame}

% ============================================================
\section{O Problema}
% ============================================================

\begin{frame}{Roraima voa às cegas no curto prazo}

  \begin{columns}[T]
    \begin{column}{0.55\textwidth}
      \begin{alertblock}{O que não existe}
        \begin{itemize}
          \item PIB estadual \textbf{trimestral} oficial do Brasil
          \item As Contas Regionais do IBGE são publicadas \textbf{apenas anualmente}
          \item A defasagem é de \textbf{até 2 anos}: os dados de 2023 foram
                publicados em outubro de 2025
          \item Roraima não tem IPCA próprio, nem PMC, PMS ou PIM-PF
        \end{itemize}
      \end{alertblock}

      \vspace{0.3cm}

      \begin{exampleblock}{O que acontece na prática}
        O gestor público toma decisões sobre orçamento, investimento e política
        econômica com base em dados que têm \textbf{2 anos de atraso}.
        É como dirigir olhando pelo espelho retrovisor.
      \end{exampleblock}
    \end{column}

    \begin{column}{0.42\textwidth}
      \begin{tikzpicture}
        \node[draw=seplanred, fill=red!8, rounded corners, thick,
              minimum width=4.2cm, minimum height=3.5cm,
              text width=4cm, align=center] (box) {%
          \color{seplanred}\Large\faExclamationTriangle\\[6pt]
          \color{seplantext}\normalsize
          \textbf{Quando o IBGE publica\\os dados de 2023...}\\[8pt]
          \large\color{seplanblue}\textbf{já estamos em\\outubro de 2025}\\[8pt]
          \normalsize\color{seplantext}
          São \textbf{22 meses}\\de defasagem
        };
      \end{tikzpicture}

      \vspace{0.4cm}
      \small\color{seplanblue}
      \textbf{Comparação:} o Banco Central publica o IBC-Br (proxy do PIB nacional)
      com \textbf{apenas 45 dias} de defasagem.
    \end{column}
  \end{columns}
\end{frame}

% ============================================================
\section{A Solução — O que é o IAET-RR}
% ============================================================

\begin{frame}{O que é o IAET-RR?}

  \begin{columns}[T]
    \begin{column}{0.52\textwidth}
      \begin{block}{Definição}
        O \textbf{IAET-RR} (Indicador de Atividade Econômica Trimestral de
        Roraima) é um índice que mede o ritmo da economia de Roraima a cada
        trimestre, em termos reais.
      \end{block}

      \vspace{0.3cm}

      \textbf{Em linguagem simples:}
      \begin{itemize}
        \item É como um \textbf{termômetro econômico} do estado
        \item Mostra se a economia cresceu, encolheu ou ficou estável
        \item Funciona trimestralmente — \textit{4 vezes por ano}
        \item Cobre todos os grandes setores: agro, indústria, comércio,
              serviços e setor público
        \item \textbf{Referência} = média de 2020 = 100
      \end{itemize}

      \vspace{0.3cm}
      \begin{exampleblock}{Metodologia}
        Equivalente ao \textbf{IBCR} do Banco Central do Brasil —
        adaptado às especificidades de Roraima.
      \end{exampleblock}
    \end{column}

    \begin{column}{0.45\textwidth}
      \centering
      \begin{tikzpicture}[scale=0.85]
        \begin{axis}[
          ybar,
          bar width=14pt,
          xlabel={Ano},
          ylabel={Variação real (\%)},
          ymin=-2, ymax=14,
          xtick={1,2,3,4,5},
          xticklabels={2021, 2022, 2023, 2024*, 2025*},
          ymajorgrids=true,
          grid style={dashed, gray!30},
          title={\small Crescimento real anual — IAET-RR},
          title style={font=\small},
          width=6.5cm, height=5.5cm,
          tick label style={font=\small},
          label style={font=\small},
          every axis plot/.style={fill opacity=0.85},
        ]
        \addplot[fill=seplanblue] coordinates {
          (1,8.19) (2,10.86) (3,4.34) (4,7.27) (5,6.92)
        };
        \addplot[fill=seplangold, opacity=0.6] coordinates {
          (1,7.23) (2,-1.23) (3,1.69) (4,4.13) (5,1.66)
        };
        \legend{\small IAET-RR, \small IBCR Norte}
        \end{axis}
      \end{tikzpicture}

      \vspace{-0.2cm}
      {\tiny * Extrapolado com tendência. Atualiza com CR-IBGE 2024 (out/2026).}
    \end{column}
  \end{columns}
\end{frame}

% ============================================================
\section{Por que Roraima Precisa Disso}
% ============================================================

\begin{frame}{Por que isso é estratégico para Roraima?}

  \begin{columns}[T]
    \begin{column}{0.48\textwidth}
      \begin{block}{\faChartLine\ \ Decisões mais informadas}
        \begin{itemize}
          \item Negociações de LDO e LOA com base em conjuntura recente
          \item Timing de investimentos públicos: saber \textit{quando} a
                economia desacelera e precisa de estímulo
          \item Monitorar o impacto de políticas setoriais em tempo real
        \end{itemize}
      \end{block}

      \vspace{0.2cm}

      \begin{block}{\faSearch\ \ Inteligência territorial}
        \begin{itemize}
          \item Identificar qual setor puxou ou freou o crescimento
          \item Comparar Roraima com o Norte e o Brasil
          \item Subsidiar notas técnicas e relatórios de conjuntura
        \end{itemize}
      \end{block}
    \end{column}

    \begin{column}{0.48\textwidth}
      \begin{block}{\faHandshake\ \ Posicionamento institucional}
        \begin{itemize}
          \item Única UF do Norte sem PIB trimestral: ser a \textbf{pioneira}
          \item Fortalecer a SEPLAN como centro de inteligência econômica
          \item Subsidiar a Assembleia Legislativa, Tribunal de Contas
                e Ministério Público
        \end{itemize}
      \end{block}

      \vspace{0.2cm}

      \begin{alertblock}{\faDollarSign\ \ O tamanho da economia em jogo}
        \centering
        PIB de RR em 2023: \textbf{R\$ 25,1 bilhões}\\
        VAB de RR em 2023: \textbf{R\$ 23,0 bilhões}\\
        {\small Crescimento nominal 2020--2023: +60\%}
      \end{alertblock}
    \end{column}
  \end{columns}

\end{frame}

% ============================================================
\section{Estrutura Setorial}
% ============================================================

\begin{frame}{O que o IAET-RR mede — 12 setores da economia}

  \begin{columns}[T]
    \begin{column}{0.50\textwidth}
      \small
      \renewcommand{\arraystretch}{1.25}
      \begin{tabular}{>{\raggedright}p{4.5cm} r}
        \toprule
        \textbf{Setor} & \textbf{\% VAB} \\
        \midrule
        \rowcolor{seplanblue!20}
        Adm. Pública (AAPP) & \textbf{46,2\%} \\
        Comércio & 12,3\% \\
        \rowcolor{seplanblue!10}
        Agropecuária & 8,9\% \\
        Imobiliário & 7,7\% \\
        \rowcolor{seplanblue!10}
        Outros Serviços & 7,6\% \\
        SIUP (energia/água) & 5,4\% \\
        \rowcolor{seplanblue!10}
        Construção & 4,9\% \\
        Financeiro & 2,8\% \\
        \rowcolor{seplanblue!10}
        Transportes & 1,9\% \\
        Ind. de Transformação & 1,3\% \\
        \rowcolor{seplanblue!10}
        Info. e Comunicação & 1,0\% \\
        Ind. Extrativas & 0,1\% \\
        \bottomrule
      \end{tabular}
      {\tiny Fonte: CR IBGE 2023. Pesos efetivos do índice usam participações de 2020.}
    \end{column}

    \begin{column}{0.47\textwidth}
      \centering
      \begin{tikzpicture}[scale=0.80]
        \pie[
          color={seplanblue, seplanblue2, seplangold, seplanblue!50,
                 seplangreen!60, seplangold!60, seplanblue!30,
                 seplanred!50, gray!60, gray!40, gray!25, gray!15},
          text=pin,
          pin edge={thin},
          radius=2.4,
          style={font=\tiny}
        ]{
          46.21/AAPP,
          12.25/Comércio,
          8.87/Agro,
          7.68/Imob.,
          7.63/Outros S.,
          5.40/SIUP,
          4.89/Constr.,
          2.78/Fin.,
          1.92/Transp.,
          1.31/Transf.,
          1.01/InfoCom,
          0.05/Extr.
        }
      \end{tikzpicture}
    \end{column}
  \end{columns}

\end{frame}

% ============================================================
\section{Metodologia}
% ============================================================

\begin{frame}{Como funciona — metodologia em linguagem simples}

  \begin{columns}[T]
    \begin{column}{0.50\textwidth}

    \textbf{Passo 1 — Coletar proxies trimestrais}\\[4pt]
    Para cada setor, usamos um \textit{termômetro} disponível em frequência mensal
    ou trimestral: emprego formal (CAGED), energia elétrica (ANEEL), passageiros
    aéreos (ANAC), folha de pagamento (SIAPE/SEPLAN), produção agrícola (IBGE).

    \vspace{0.4cm}

    \textbf{Passo 2 — Calibrar pelo IBGE (Denton-Cholette)}\\[4pt]
    As proxies são ajustadas para que a \textbf{média anual} de cada setor bata
    exatamente com o número oficial do IBGE (Contas Regionais). Garantia de
    credibilidade e ancoragem institucional.

    \vspace{0.4cm}

    \textbf{Passo 3 — Agregar}\\[4pt]
    Os setores são combinados com pesos proporcionais ao tamanho de cada um
    na economia de RR (fórmula de Laspeyres, padrão das Contas Nacionais).

    \end{column}

    \begin{column}{0.47\textwidth}
      \centering
      \begin{tikzpicture}[
        node distance=0.45cm,
        box/.style={draw=seplanblue, fill=seplanlightblue,
                    rounded corners=4pt, text width=5cm,
                    minimum height=0.75cm, align=center,
                    font=\small},
        arrow/.style={-Stealth, thick, seplanblue},
        label/.style={font=\footnotesize, text=seplanblue, align=center}
      ]
        \node[box, fill=seplangold!25, draw=seplangold] (s1)
          {\textbf{10+ fontes de dados}\\{\tiny CAGED, ANEEL, ANAC, SIAPE, IBGE\ldots}};
        \node[box, below=of s1] (s2)
          {\textbf{Proxies trimestrais por setor}\\{\tiny Índices de volume: 2020T1 a 2025T4}};
        \node[box, below=of s2] (s3)
          {\textbf{Calibração Denton-Cholette}\\{\tiny Ancora ao VAB real das CR IBGE}};
        \node[box, below=of s3] (s4)
          {\textbf{Agregação Laspeyres}\\{\tiny Pesos = participação no VAB de 2020}};
        \node[box, below=of s4, fill=seplangreen!15, draw=seplangreen] (s5)
          {\textbf{IAET-RR — Índice trimestral}\\{\tiny Base 2020 = 100 $\cdot$ NSA e SA}};

        \draw[arrow] (s1) -- (s2);
        \draw[arrow] (s2) -- (s3);
        \draw[arrow] (s3) -- (s4);
        \draw[arrow] (s4) -- (s5);

        \node[label, right=0.1cm of s3, text width=1.8cm]
          {\tiny Garante\\ coerência\\ com IBGE};
        \node[label, right=0.1cm of s4, text width=1.8cm]
          {\tiny Padrão\\ Contas\\ Nacionais};
      \end{tikzpicture}
    \end{column}
  \end{columns}

\end{frame}

\begin{frame}{Ancoragem ao IBGE — o que isso garante}

  \begin{columns}[T]
    \begin{column}{0.50\textwidth}
      \begin{block}{O que é o Denton-Cholette?}
        É uma técnica matemática que \textbf{interpola} a série anual do IBGE
        dentro de cada ano, usando a evolução mensal das nossas proxies como guia.

        \vspace{0.3cm}
        Resultado: o IAET-RR nunca diverge dos números oficiais do IBGE nos anos
        em que há Contas Regionais disponíveis.
      \end{block}

      \vspace{0.3cm}

      \begin{exampleblock}{Verificação técnica}
        Ancoragem \textbf{perfeita} em todos os 4 anos com dados:\\
        \begin{tabular}{llr}
          2020 & média trimestral & 100,000 \\
          2021 & média trimestral & 108,192 \\
          2022 & média trimestral & 119,945 \\
          2023 & média trimestral & 125,152 \\
        \end{tabular}\\
        Desvio em relação ao IBGE: $< 0{,}001$
      \end{exampleblock}
    \end{column}

    \begin{column}{0.47\textwidth}
      \begin{alertblock}{Por que usamos volume, não valor?}
        O indicador mede \textbf{volume físico de atividade} — não valor em
        reais. Isso é essencial porque:

        \begin{itemize}
          \item Remove o efeito da inflação
          \item Permite comparar entre trimestres e anos
          \item É o mesmo padrão do IBCR do Banco Central
        \end{itemize}

        \vspace{0.3cm}
        \textbf{Decisão metodológica de 2026:} o benchmark anual passou a usar
        o \textit{índice de volume} das Contas Regionais (e não o VAB nominal).
        Resultado: crescimentos saíram de +20\% nominal para +4\% real em 2023
        — tecnicamente correto e comparável ao IBCR.
      \end{alertblock}
    \end{column}
  \end{columns}

\end{frame}

% ============================================================
\section{Fontes de Dados e Automação}
% ============================================================

\begin{frame}{10+ fontes integradas — coleta automatizada}

  \begin{columns}[T]
    \begin{column}{0.50\textwidth}
      \small
      \renewcommand{\arraystretch}{1.2}
      \begin{tabular}{>{\raggedright}p{2.8cm} >{\raggedright}p{2.8cm}}
        \toprule
        \textbf{Fonte} & \textbf{Dado} \\
        \midrule
        IBGE / SIDRA & Produção agropecuária \\
        \rowcolor{seplanblue!10}
        SIAPE / MF & Folha federal por UF \\
        SEPLAN-RR & Folha estadual \\
        \rowcolor{seplanblue!10}
        SICONFI / STN & Folha municipal \\
        ANEEL (SAMP) & Energia elétrica por classe \\
        \rowcolor{seplanblue!10}
        MTE / CAGED & Emprego formal (7 seções) \\
        ANAC & Passageiros/carga aéreos \\
        \rowcolor{seplanblue!10}
        ANP & Vendas de diesel \\
        BCB (Estban/SCR) & Crédito e depósitos \\
        \rowcolor{seplanblue!10}
        IBGE (IPCA) & Deflator nacional \\
        CR IBGE & Benchmark anual (volume) \\
        \bottomrule
      \end{tabular}
    \end{column}

    \begin{column}{0.47\textwidth}
      \begin{block}{\faRobot\ \ Pipeline automatizado em R}
        \begin{itemize}
          \item \textbf{8 scripts} encadeados (00 a 05f)
          \item Atualização trimestral com um único comando
          \item Download automático da maioria das fontes via API
          \item Validação automática em 4 eixos
          \item Exportação automática para \texttt{.csv} e \texttt{.xlsx}
        \end{itemize}
      \end{block}

      \vspace{0.3cm}

      \begin{exampleblock}{\faClock\ \ Tempo de atualização}
        \centering
        \textbf{$\approx$ 30 minutos} por trimestre\\
        (após a publicação dos dados de fonte)
      \end{exampleblock}

      \vspace{0.2cm}
      {\small Código versionado em repositório GitHub público
      da SEPLAN/RR — rastreabilidade total.}
    \end{column}
  \end{columns}

\end{frame}

% ============================================================
\section{Resultados}
% ============================================================

\begin{frame}{Resultados — Série trimestral 2020--2025}

  \begin{columns}[T]
    \begin{column}{0.58\textwidth}
      \centering
      \begin{tikzpicture}
        \begin{axis}[
          width=8.5cm, height=6cm,
          xlabel={Trimestre},
          ylabel={Índice (2020 = 100)},
          xtick={1,5,9,13,17,21},
          xticklabels={2020T1,2021T1,2022T1,2023T1,2024T1,2025T1},
          x tick label style={rotate=30, anchor=east, font=\tiny},
          ymajorgrids=true,
          grid style={dashed, gray!25},
          title={\small IAET-RR — Índice Geral (base 2020=100)},
          title style={font=\small},
          tick label style={font=\tiny},
          label style={font=\small},
          legend style={font=\tiny, at={(0.02,0.98)}, anchor=north west},
          ymin=75, ymax=160,
        ]
        % Índice geral NSA
        \addplot[seplanblue, very thick, mark=*, mark size=1.5pt] coordinates {
          (1,100.95)(2,88.06)(3,101.59)(4,109.40)
          (5,103.83)(6,96.17)(7,116.27)(8,116.50)
          (9,109.78)(10,111.89)(11,129.49)(12,128.62)
          (13,117.52)(14,117.38)(15,134.89)(16,130.83)
          (17,128.36)(18,131.98)(19,135.93)(20,140.74)
          (21,140.58)(22,142.55)(23,144.71)(24,146.37)
        };
        % AAPP
        \addplot[seplangold, dashed, thick] coordinates {
          (1,98.66)(2,97.86)(3,98.97)(4,104.51)
          (5,107.10)(6,78.63)(7,112.41)(8,114.63)
          (9,107.65)(10,106.90)(11,106.47)(12,108.76)
          (13,99.94)(14,105.71)(15,114.36)(16,119.95)
          (17,120.65)(18,121.36)(19,122.07)(20,122.78)
          (21,123.50)(22,124.22)(23,124.95)(24,125.68)
        };
        \addlegendentry{Índice Geral}
        \addlegendentry{Adm. Pública}
        % COVID annotation
        \draw[seplanred, thick, ->] (axis cs:2,78) -- (axis cs:2,89)
          node[above, font=\tiny, text=seplanred] {COVID-19\\T2/2020};
        \end{axis}
      \end{tikzpicture}
    \end{column}

    \begin{column}{0.38\textwidth}
      \small
      \begin{block}{Leituras-chave}
        \renewcommand{\arraystretch}{1.3}
        \begin{tabular}{lr}
          \textbf{Período} & \textbf{Variação} \\
          \hline
          2020 (COVID) & \color{seplanred}$-$9,8\% T1→T2 \\
          Recuperação T3 & \color{seplangreen}+15,4\% \\
          \hline
          2021 (anual) & \color{seplangreen}+8,2\% \\
          2022 (anual) & \color{seplangreen}+10,9\% \\
          2023 (anual) & \color{seplangreen}+4,3\% \\
          \hline
          2024* (anual) & \color{seplangreen}+7,3\% \\
          2025* (anual) & \color{seplangreen}+6,9\% \\
        \end{tabular}
      \end{block}

      \vspace{0.2cm}
      {\small \color{seplanblue}
      \textbf{Pico observado:} 2023T3\\
      Índice = \textbf{134,9}\\
      (puxado pela soja e SIUP)
      }

      \vspace{0.2cm}
      {\tiny * Extrapolado com tendência geométrica 2022--2023.
      Atualiza com CR IBGE 2024 (out/2026).}
    \end{column}
  \end{columns}

\end{frame}

\begin{frame}{Quem puxou o crescimento — análise setorial}

  \begin{columns}[T]
    \begin{column}{0.55\textwidth}
      \centering
      \begin{tikzpicture}
        \begin{axis}[
          ybar,
          bar width=9pt,
          width=8cm, height=5.8cm,
          xlabel={},
          ylabel={Índice (2020=100)},
          xtick={1,2,3,4},
          xticklabels={Agro,AAPP,Indústria,Serviços},
          x tick label style={font=\small},
          ymajorgrids=true,
          grid style={dashed, gray!25},
          title={\small Média anual por bloco setorial},
          title style={font=\small},
          ymin=0, ymax=420,
          legend style={font=\tiny, at={(0.02,0.98)}, anchor=north west},
          tick label style={font=\tiny},
          label style={font=\small},
        ]
        \addplot[fill=seplanblue!70] coordinates {(1,100)(2,100)(3,100)(4,100)};
        \addplot[fill=seplanblue] coordinates {(1,150.8)(2,108.2)(3,112.1)(4,115.1)};
        \addplot[fill=seplangold] coordinates {(1,270.2)(2,110.5)(3,133.5)(4,120.5)};
        \addplot[fill=seplangold!70] coordinates {(1,312.7)(2,110.0)(3,143.9)(4,127.6)};
        \legend{2020,2021,2022,2023}
        \end{axis}
      \end{tikzpicture}
    \end{column}

    \begin{column}{0.42\textwidth}
      \begin{block}{Interpretação}
        \textbf{Agropecuária} foi o setor de maior expansão relativa no período
        (índice 2023 $\approx$ 313 — mais que triplicou o nível de 2020).
        Reflecte o boom da soja roraimense.

        \vspace{0.25cm}

        \textbf{Indústria} cresceu sustentadamente — puxada pela expansão do
        SIUP (energia/água) e obras de infraestrutura (Construção).

        \vspace{0.25cm}

        \textbf{AAPP} cresceu moderadamente (+10\% em 3 anos), condizente
        com crescimento real da folha pública.

        \vspace{0.25cm}

        \textbf{Serviços privados} cresceram +28\% acumulado, refletindo
        a expansão do consumo privado.
      \end{block}
    \end{column}
  \end{columns}

\end{frame}

\begin{frame}{IAET-RR vs. referências nacionais}

  \begin{columns}[T]
    \begin{column}{0.50\textwidth}
      \centering
      \begin{tikzpicture}
        \begin{axis}[
          ybar=3pt,
          bar width=11pt,
          width=7.5cm, height=5.5cm,
          xlabel={Ano},
          ylabel={Variação anual (\%)},
          xtick={1,2,3,4,5},
          xticklabels={2021,2022,2023,2024,2025},
          ymajorgrids=true,
          grid style={dashed, gray!25},
          ymin=-4, ymax=14,
          title={\small Variação anual real (\%)},
          title style={font=\small},
          legend style={font=\tiny, at={(0.98,0.98)}, anchor=north east},
          tick label style={font=\tiny},
          label style={font=\small},
        ]
        \addplot[fill=seplanblue] coordinates {
          (1,8.19)(2,10.86)(3,4.34)(4,7.27)(5,6.92)};
        \addplot[fill=seplangold] coordinates {
          (1,7.23)(2,-1.23)(3,1.69)(4,4.13)(5,1.66)};
        \addplot[fill=seplangreen!60] coordinates {
          (1,4.2)(2,2.8)(3,2.7)(4,3.7)(5,2.5)};
        \legend{IAET-RR, IBCR Norte, IBC-Br}
        \end{axis}
      \end{tikzpicture}
    \end{column}

    \begin{column}{0.47\textwidth}
      \begin{block}{O que os números dizem}
        \begin{itemize}
          \item Roraima cresceu \textbf{acima do Norte} em todos os anos
                com dados disponíveis
          \item Em 2022, enquanto o Norte encolheu (--1,2\%), Roraima
                avançou \textbf{+10,9\%}
          \item O crescimento de RR é estruturalmente diferente do Norte
                (dominado por Pará e Amazonas)
          \item RR se destaca no contexto regional — dado que antes era
                \textbf{invisível} sem este indicador
        \end{itemize}
      \end{block}

      \vspace{0.2cm}

      \begin{exampleblock}{Correlação com IBCR}
        Correlação do IAET-RR com o IBC-Br em nível: \textbf{0,934}\\
        Alta correlação confirma a qualidade metodológica do índice.
      \end{exampleblock}
    \end{column}
  \end{columns}

\end{frame}

\begin{frame}{Ajuste sazonal — lendo o ciclo sem ruído}

  \begin{columns}[T]
    \begin{column}{0.55\textwidth}
      \centering
      \begin{tikzpicture}
        \begin{axis}[
          width=8cm, height=5.5cm,
          xlabel={Trimestre},
          ylabel={Índice (2020=100)},
          xtick={1,5,9,13},
          xticklabels={2020T1,2021T1,2022T1,2023T1},
          x tick label style={rotate=25, anchor=east, font=\tiny},
          ymajorgrids=true,
          grid style={dashed, gray!25},
          title={\small NSA vs. Dessazonalizado (SA) — 2020--2023},
          title style={font=\small},
          tick label style={font=\tiny},
          label style={font=\small},
          legend style={font=\tiny, at={(0.02,0.98)}, anchor=north west},
          ymin=80, ymax=145,
        ]
        % NSA
        \addplot[seplanblue, thick, mark=*, mark size=1pt] coordinates {
          (1,100.95)(2,88.06)(3,101.59)(4,109.40)
          (5,103.83)(6,96.17)(7,116.27)(8,116.50)
          (9,109.78)(10,111.89)(11,129.49)(12,128.62)
          (13,117.52)(14,117.38)(15,134.89)(16,130.83)
        };
        % SA (NSA - fator sazonal médio aproximado)
        \addplot[seplanred, thick, dashed, mark=square*, mark size=1pt] coordinates {
          (1,103.62)(2,96.84)(3,95.54)(4,103.80)
          (5,107.08)(6,104.43)(7,109.91)(8,111.53)
          (9,113.02)(10,120.18)(11,123.13)(12,123.54)
          (13,120.82)(14,126.10)(15,128.53)(16,125.87)
        };
        \legend{Sem ajuste (NSA),Dessazonalizado (SA)}
        \end{axis}
      \end{tikzpicture}
    \end{column}

    \begin{column}{0.42\textwidth}
      \begin{block}{O que é ajuste sazonal?}
        Em Roraima, o T3 (julho--setembro) é naturalmente mais alto por
        causa da \textbf{colheita da soja}. O ajuste sazonal retira esse
        efeito, revelando a tendência real da economia.
      \end{block}

      \vspace{0.3cm}

      \begin{exampleblock}{Método: X-13ARIMA-SEATS}
        O mesmo algoritmo usado pelo \textbf{Bureau of Census dos EUA} e
        pelo \textbf{Eurostat}. É o padrão ouro mundial para ajuste
        sazonal de séries econômicas.

        \vspace{0.2cm}
        Amplitude sazonal do índice geral:\\
        \textbf{$\approx$ 15 pontos} (pico T3 vs. vale T2)
      \end{exampleblock}

      \vspace{0.2cm}
      {\small Com a série dessazonalizada, é possível comparar qualquer
      trimestre com qualquer outro sem distorção de calendário.}
    \end{column}
  \end{columns}

\end{frame}

% ============================================================
\section{Produto — Séries Disponíveis}
% ============================================================

\begin{frame}{O que está disponível hoje}

  \begin{columns}[T]
    \begin{column}{0.50\textwidth}
      \begin{block}{\faTable\ \ Séries geradas (2020T1--2025T4)}
        \renewcommand{\arraystretch}{1.2}
        \small
        \begin{tabular}{>{\raggedright}p{3.5cm} l}
          \textbf{Arquivo} & \textbf{Conteúdo} \\
          \hline
          \texttt{indice\_geral\_rr.csv} & Índice geral + setoriais \\
          \texttt{indice\_geral\_rr\_sa.csv} & Série dessazonalizada \\
          \texttt{fatores\_sazonais.csv} & Fatores X-13 por setor \\
          \texttt{vab\_nominal\_rr\_reais.csv} & VAB nominal em R\$ mi \\
          \texttt{indice\_nominal\_rr.csv} & Índice nominal (c/ deflator) \\
          \texttt{IAET\_RR\_series.xlsx} & Excel completo (6 abas) \\
        \end{tabular}
      \end{block}

      \vspace{0.3cm}

      \begin{exampleblock}{VAB nominal trimestral (R\$ milhões)}
        \small
        \begin{tabular}{lr}
          2020 (soma anual) & R\$ 14.531 mi \\
          2021 & R\$ 16.331 mi \\
          2022 & R\$ 19.301 mi \\
          2023 & R\$ 23.297 mi \\
          2024* & R\$ 26.673 mi \\
          2025* & R\$ 29.944 mi \\
        \end{tabular}\\
        {\tiny * Extrapolado.}
      \end{exampleblock}
    \end{column}

    \begin{column}{0.47\textwidth}
      \centering
      \begin{tikzpicture}
        \begin{axis}[
          ybar,
          bar width=16pt,
          width=6.5cm, height=5.5cm,
          xlabel={Ano},
          ylabel={R\$ bilhões},
          xtick={1,2,3,4,5,6},
          xticklabels={2020,2021,2022,2023,2024*,2025*},
          ymajorgrids=true,
          grid style={dashed, gray!25},
          ymin=0, ymax=35,
          title={\small VAB Nominal de Roraima (R\$ bi)},
          title style={font=\small},
          tick label style={font=\tiny},
          label style={font=\small},
          nodes near coords,
          nodes near coords align={vertical},
          every node near coord/.style={font=\tiny, color=seplantext},
        ]
        \addplot[fill=seplanblue, opacity=0.85] coordinates {
          (1,14.5)(2,16.3)(3,19.3)(4,23.3)(5,26.7)(6,29.9)
        };
        \end{axis}
      \end{tikzpicture}

      \vspace{0.2cm}
      {\small\color{seplanblue}
      \textbf{+106\%} de crescimento nominal\\em 5 anos (2020--2025*)}
    \end{column}
  \end{columns}

\end{frame}

% ============================================================
\section{Status e Próximos Passos}
% ============================================================

\begin{frame}{Estado atual e o que vem a seguir}

  \begin{columns}[T]
    \begin{column}{0.48\textwidth}
      \begin{exampleblock}{\faCheckCircle\ \ O que está pronto}
        \begin{itemize}
          \item[\color{seplangreen}\faCheck] Pipeline de produção completo (8 scripts)
          \item[\color{seplangreen}\faCheck] Séries 2020T1--2025T4 geradas
          \item[\color{seplangreen}\faCheck] Ajuste sazonal (X-13ARIMA-SEATS)
          \item[\color{seplangreen}\faCheck] VAB nominal trimestral em R\$ mi
          \item[\color{seplangreen}\faCheck] Validação em 4 eixos automatizada
          \item[\color{seplangreen}\faCheck] Excel de publicação (6 abas)
          \item[\color{seplangreen}\faCheck] Código público no GitHub
          \item[\color{seplangreen}\faCheck] Documentação metodológica completa
        \end{itemize}
      \end{exampleblock}
    \end{column}

    \begin{column}{0.48\textwidth}
      \begin{alertblock}{\faTools\ \ Próximas entregas}
        \begin{itemize}
          \item[\color{seplanred}\faClock] \textbf{Dashboard interativo} (Shiny)\\
                {\small Gráficos, tabelas e download de dados}

          \vspace{0.1cm}
          \item[\color{seplanred}\faClock] \textbf{Nota técnica} (Quarto/PDF)\\
                {\small Metodologia + análise conjuntural}

          \vspace{0.1cm}
          \item[\color{seplangold}\faClock] \textbf{CR IBGE 2024} (out/2026)\\
                {\small Substituir extrapolações 2024 por dados reais}

          \vspace{0.1cm}
          \item[\color{seplangold}\faClock] \textbf{ICMS SEFAZ-RR}\\
                {\small Melhorar proxy do Comércio e gerar PIB nominal}
        \end{itemize}
      \end{alertblock}
    \end{column}
  \end{columns}

  \vspace{0.3cm}

  \begin{block}{\faLightbulb\ \ Compromisso de atualização}
    O IAET-RR será atualizado \textbf{trimestralmente}, com publicação nos
    primeiros 45 dias após o encerramento de cada trimestre. Tempo de
    processamento: aproximadamente 30 minutos por ciclo.
  \end{block}

\end{frame}

% ============================================================
\section{Usos e Aplicações}
% ============================================================

\begin{frame}{Para quê serve na prática — usos imediatos}

  \begin{columns}[T]
    \begin{column}{0.48\textwidth}

      \begin{block}{\faBuilding\ \ Para a SEPLAN/RR}
        \begin{itemize}
          \item Embasar projeções de receita e despesa na LDO/LOA
          \item Identificar setores em desaceleração antes de decisões
                de investimento público
          \item Publicar boletim econômico trimestral com dados próprios
          \item Subsidiar o Conselho de Desenvolvimento Econômico do Estado
        \end{itemize}
      \end{block}

      \vspace{0.25cm}

      \begin{block}{\faGavel\ \ Para outros órgãos}
        \begin{itemize}
          \item \textbf{SEFAZ-RR}: projetar arrecadação tributária
          \item \textbf{ALES}: acompanhar execução orçamentária
          \item \textbf{TCE-RR}: auditoria com contexto econômico
          \item \textbf{Setor privado}: decisões de investimento e expansão
        \end{itemize}
      \end{block}
    \end{column}

    \begin{column}{0.48\textwidth}
      \begin{alertblock}{\faGlobe\ \ Posicionamento estratégico}
        \textbf{Somos a única UF do Norte} a construir um indicador
        trimestral com essa cobertura setorial.

        \vspace{0.3cm}
        Outros estados que já têm indicadores próprios:
        \begin{itemize}
          \item São Paulo (SEADE — IPR-SP)
          \item Rio Grande do Sul (FEE — ICV-RS)
          \item Minas Gerais (FIEMG — IEA-MG)
          \item Ceará (IPECE — IPCA-CE trimestral)
        \end{itemize}

        \vspace{0.3cm}
        \textbf{Com o IAET-RR, Roraima entra nesse seleto grupo.}
      \end{alertblock}

      \vspace{0.2cm}
      {\small\color{seplanblue}
      O indicador pode fundamentar publicações em periódicos acadêmicos
      e parcerias com UFAM, UFRR e IPAM.}
    \end{column}
  \end{columns}

\end{frame}

% ============================================================
\section{Encerramento}
% ============================================================

\begin{frame}{O que pedimos — e o que entregamos}

  \begin{columns}[T]
    \begin{column}{0.50\textwidth}
      \begin{block}{O que a DIEAS entregou}
        \begin{itemize}
          \item Metodologia original, documentada e auditável
          \item Pipeline automatizado e reprodutível
          \item 24 trimestres de série histórica (2020--2025)
          \item Abertura por 4 blocos setoriais + 12 atividades
          \item Série dessazonalizada e deflator implícito
          \item VAB nominal trimestral em R\$ milhões
          \item Código público — transparência total
        \end{itemize}
      \end{block}
    \end{column}

    \begin{column}{0.47\textwidth}
      \begin{exampleblock}{O que pedimos}
        \begin{enumerate}
          \item \textbf{Aval da Coordenação da CGEES} para avançar para
                publicação oficial (dashboard + nota técnica)

          \vspace{0.2cm}
          \item \textbf{Apoio do Secretário} para obter o ICMS da
                SEFAZ-RR (necessário para melhorar o bloco Comércio
                e gerar o PIB nominal)

          \vspace{0.2cm}
          \item \textbf{Publicação do IAET-RR} como produto oficial
                da SEPLAN/RR no portal do Estado
        \end{enumerate}
      \end{exampleblock}
    \end{column}
  \end{columns}

  \vspace{0.4cm}

  \begin{center}
    \large\color{seplanblue}
    \textbf{Roraima tem economia. Agora tem como medi-la.}
  \end{center}

\end{frame}

% --- Slide final --------------------------------------------
{
\setbeamertemplate{footline}{}
\begin{frame}[plain]
  \begin{center}
    \vspace{0.5cm}
    {\Large\color{seplanblue}\textbf{Obrigado}}\\[0.5cm]
    {\color{seplanblue}\rule{8cm}{0.5pt}}\\[0.5cm]
    {\normalsize
      \textbf{Yuri Cesar de Lima e Silva}\\
      Chefe da DIEAS — SEPLAN/RR\\[0.3cm]
      Coordenação-Geral de Estudos Econômicos e Sociais\\
      Secretaria de Estado do Planejamento e Desenvolvimento de Roraima\\[0.5cm]
    }
    {\small\color{gray}
      Código-fonte e documentação:\\
      \texttt{github.com/yuricesarsilva/painel\_pib\_trimestral}\\[0.3cm]
      Documentação técnica:\\
      \texttt{plano\_projeto.md} $\cdot$
      \texttt{plano\_reforma\_indicador\_real.md} $\cdot$
      \texttt{checklist.md}
    }\\[0.5cm]
    {\color{seplanblue}\rule{8cm}{0.5pt}}\\[0.3cm]
    {\footnotesize\color{gray} Abril de 2026}
  \end{center}
\end{frame}
}

% --- Apêndice -----------------------------------------------
\appendix

\begin{frame}{Apêndice — Série trimestral completa}
  \small
  \renewcommand{\arraystretch}{1.1}
  \begin{tabular}{lccccc}
    \toprule
    \textbf{Período} & \textbf{Geral} & \textbf{Agro} & \textbf{AAPP} & \textbf{Indústria} & \textbf{Serviços} \\
    \midrule
    2020T1 & 100,9 &  31,1 &  98,7 & 100,6 & 104,9 \\
    2020T2 &  88,1 &  23,1 &  97,9 &  91,5 &  87,6 \\
    2020T3 & 101,6 & 215,2 &  99,0 &  92,7 &  97,5 \\
    2020T4 & 109,4 & 130,7 & 104,5 & 115,2 & 109,9 \\
    \hline
    2021T1 & 103,8 &  29,1 & 107,1 & 113,9 & 106,7 \\
    2021T2 &  96,2 &  20,0 &  78,6 & 105,7 & 106,5 \\
    2021T3 & 116,3 & 283,8 & 112,4 & 109,4 & 111,9 \\
    2021T4 & 116,5 & 166,4 & 114,6 & 119,5 & 117,0 \\
    \hline
    2022T1 & 109,8 &  25,1 & 107,6 & 129,4 & 115,9 \\
    2022T2 & 111,9 &  18,6 & 106,9 & 124,0 & 120,3 \\
    2022T3 & 129,5 & 382,3 & 106,5 & 134,1 & 127,5 \\
    2022T4 & 128,6 & 213,2 & 108,8 & 146,1 & 132,6 \\
    \hline
    2023T1 & 117,5 &  25,9 &  99,9 & 137,1 & 128,7 \\
    2023T2 & 117,4 &  18,5 & 105,7 & 141,5 & 126,3 \\
    2023T3 & 134,9 & 456,9 & 114,4 & 144,9 & 126,7 \\
    2023T4 & 130,8 & 249,6 & 120,0 & 151,9 & 128,6 \\
    \bottomrule
  \end{tabular}\\[4pt]
  {\tiny Base 2020 = 100. Dados reais (ancoragem Denton-Cholette ao volume CR IBGE). 2024--2025 extrapolados — ver CSV completo.}
\end{frame}

\begin{frame}{Apêndice — Qualidade das proxies por setor}
  \small
  \renewcommand{\arraystretch}{1.2}
  \begin{tabular}{>{\raggedright}p{3.5cm} >{\raggedright}p{2.2cm} >{\raggedright}p{2.0cm} l}
    \toprule
    \textbf{Setor} & \textbf{Proxy principal} & \textbf{Tipo} & \textbf{Qualidade} \\
    \midrule
    Agropecuária & PAM + PPM & Físico & \color{seplangreen}\textbf{Forte} \\
    Adm. Pública & Folha (SIAPE+SEPLAN) & Valor real & \color{seplangreen}\textbf{Forte} \\
    SIUP & Energia ANEEL & Físico & \color{seplangreen}\textbf{Forte} \\
    Comércio & Energia + CAGED G & Misto & \color{seplangold}\textbf{Médio}* \\
    Transportes & ANAC + ANP & Físico & \color{seplangreen}\textbf{Forte} \\
    Financeiro & BCB SCR + Estban & Deflacionado & \color{seplangold}\textbf{Médio} \\
    Construção & CAGED F & Insumo & \color{seplangold}\textbf{Aceitável} \\
    Outros Serv. & CAGED I+M+N+P+Q & Insumo & \color{seplangold}\textbf{Aceitável} \\
    Info. e Com. & CAGED J & Insumo & \color{seplangold}\textbf{Aceitável} \\
    Imobiliário & Interpolação CR & Tendência & \color{seplanred}\textbf{Fraca}** \\
    Extrativas & Interpolação CR & Tendência & \color{seplanred}\textbf{Fraca}** \\
    \bottomrule
  \end{tabular}\\[4pt]
  {\tiny * Melhora com ICMS da SEFAZ-RR.
   ** Peso combinado $<$ 8\% — impacto limitado no índice geral.}
\end{frame}

\end{document}
